% ── Korelační analýza: pokrytí KS a ukazatele trhu práce ─────────────────────

\paragraph{Datová základna a pokrytí zemí EU27}\label{par:korelace_data}

Korelační analýza pracuje s~párovým panelem (\textit{geo\,$\times$\,rok}) sestaveným z~následujících zdrojů.
Pokrytí \ac{KS} pochází z~databáze OECD/AIAS ICTWSS~\cite{oecd_aias_ictwss_CBC_ERB_pct}: pro 22~zemí EU27 je použita upravená míra \textit{AdjCov} zachycující podíl zaměstnanců pokrytých jakoukoliv kolektivní smlouvou bez ohledu na odborové členství; pro Německo~(DE) a~Slovensko~(SK) je \textit{AdjCov} nahrazena mírou ERB z~průzkumu OECD/CBC, neboť \textit{AdjCov} pro tyto země stagnuje přibližně na hodnotách z~roku 1990 resp.~2015.
Hustota odborů pochází z~téže databáze (ukazatel TUD~\cite{oecd_aias_ictwss_TUD_pct}).
Výstupní proměnné pocházejí z~Eurostatu: průměrná týdenní pracovní doba (lfsa\_ewhun2~\cite{eurostat_lfsa_ewhun2_HR_weekly}), HDP~na~obyvatele v~PPS (nama\_10\_pc~\cite{eurostat_nama_10_pc_PPS_EU27eq100}), Giniho koeficient disponibilního příjmu (ilc\_di12~\cite{eurostat_ilc_di12_Gini}) a~čistý roční příjem v~PPS (earn\_nt\_net~\cite{eurostat_earn_nt_net_PPS_AW100}).
Datová řada začíná rokem 2004 shodně ve všech souborech.

Ze 27~členských zemí EU jsou do analýzy zahrnuty \textbf{čtyřiadvacet}; z~dat jsou vyjmuty tři:
\begin{itemize}
  \item \textbf{Irsko~(IE)} je záměrně vyřazeno jako statistický outlier.
    Jeho HDP/ob.~v~PPS dosahuje přibližně 237~(EU27\,=\,100), přičemž tato hodnota je způsobena přesunem zisků nadnárodních technologických a~farmaceutických korporací se sídlem v~Irsku pro účely daňové optimalizace v~prostoru EU (transfer pricing, relokace duševního vlastnictví).
    Skutečná životní úroveň, vyjádřená upravenými národními příjmy GNI* (CSO Ireland), se pohybuje okolo indexu~135.
    Zahrnutí Irska by systematicky zkreslilo regresní koeficient $\hat{\beta}$ u~páru pokrytí~vs.~HDP a~tím deformovalo predikci pro zbývající země.
  \item \textbf{Estonsko~(EE) a~Lotyšsko~(LV)} --- databáze ICTWSS \textit{AdjCov} pro tyto země neobsahuje žádné záznamy; jde o~strukturálně chybějící hodnoty (MCAR), jejichž doplnění bez externího zdroje není opodstatněné.
  \item \textbf{Lucembursko~(LU)} --- data \textit{AdjCov} jsou dostupná pouze před rokem 2004; po aplikaci filtru \texttt{START\_YEAR\,=\,2004} vypadávají z~panelu.
\end{itemize}

Zbývajících 24~zemí pokrývá všechny geografické a~institucionální skupiny EU27: nordické, kontinentální, středomořské, středoevropské a~balkánské (s~výjimkou EE a~LV).
Panel není vyvážený: počet párových pozorování~$n$ se liší dle dostupnosti výstupní proměnné (viz tabulku~\ref{tab:coverage_correlation_table}).
Nejřídší je pár pokrytí~vs.~Gini ($n = 163$), neboť šetření EU-SILC má mezery zejména pro roky 2020--2021 (omezený sběr dat v~době pandemie COVID-19), a~Gini pro některé nové členské státy chybí před rokem 2007.
Nejhustší panel má pár pokrytí~vs.~HDP ($n = 299$) díky úplné dostupnosti Eurostat / nama\_10\_pc.

\paragraph{Matematický aparát}\label{par:korelace_math}

Jsou použity tři statistické míry asociace a~jeden predikční model.

\emph{Pearsonův korelační koeficient} měří sílu a~směr lineárního vztahu mezi proměnnými~$x$ a~$y$:
\begin{equation}\label{eq:pearson_r}
  r_{xy}
  = \frac{\displaystyle\sum_{i=1}^{n}(x_i - \bar{x})\,(y_i - \bar{y})}
         {\displaystyle\sqrt{\sum_{i=1}^{n}(x_i-\bar{x})^2
                             \cdot \sum_{i=1}^{n}(y_i-\bar{y})^2}}
  \;\in\; [-1,\,1].
\end{equation}

\emph{Spearmanův koeficient~$\rho$} je Pearsonovo~$r$ aplikované na hodnosti (ranky) pozorování namísto původních hodnot:
\begin{equation}\label{eq:spearman_rho}
  \rho = r\!\bigl(R(x),\,R(y)\bigr)
       = 1 - \frac{6\,\displaystyle\sum_{i=1}^{n} d_i^{\,2}}{n\,(n^2-1)},
  \quad d_i = R(x_i) - R(y_i),
\end{equation}
kde pravý výraz platí při absenci shodných hodnot (\textit{ties}); obecně platí levý tvar.
Výrazná divergence~$|r - \rho|$ signalizuje přítomnost vlivných odlehlých pozorování, která posouvají lineární vztah.

\emph{Prostá lineární regrese} metodou nejmenších čtverců (OLS):
\begin{equation}\label{eq:ols}
  \hat{y} = \hat{\alpha} + \hat{\beta}\,x,
  \qquad
  \hat{\beta}
    = \frac{\displaystyle\sum_{i=1}^{n}(x_i-\bar{x})\,(y_i-\bar{y})}
           {\displaystyle\sum_{i=1}^{n}(x_i-\bar{x})^2},
  \quad
  \hat{\alpha} = \bar{y} - \hat{\beta}\,\bar{x}.
\end{equation}
Parametr~$\hat{\beta}$ udává průměrnou změnu $\hat{y}$ při nárůstu pokrytí \ac{KS} o~jeden procentní bod.

\emph{Koeficient determinace}~$R^2$ vyjadřuje podíl rozptylu~$y$ vysvětlený regresním modelem:
\begin{equation}\label{eq:r2}
  R^2
  = 1 - \frac{\displaystyle\sum_{i=1}^{n}(y_i - \hat{y}_i)^2}
             {\displaystyle\sum_{i=1}^{n}(y_i - \bar{y})^2}
  = r^2 \;\in\; [0,\,1],
\end{equation}
kde rovnost $R^2 = r^2$ platí pro prostou regresi s~konstantou.

\emph{Predikce dopadu změny pokrytí}~$\Delta x$ na výstupní proměnnou:
\begin{equation}\label{eq:delta_pred}
  \Delta\hat{y} = \hat{\beta}\,\Delta x.
\end{equation}
Tento vztah je využit v~modelovém výpočtu pro ČR (viz~odst.~\ref{par:model_cz}).
Síla korelací je hodnocena dle Evansovy (1996) škály: $|r| < 0{,}20$ velmi slabá;
$0{,}20$--$0{,}39$ slabá; $0{,}40$--$0{,}59$ střední; $0{,}60$--$0{,}79$ silná;
$\geq 0{,}80$ velmi silná.

\paragraph{Kombinovaný pohled na asociace pokrytí KS s trhem práce}\label{par:scatter_combined}

Kombinovaný rozptylový diagram (obr.~\ref{fig:scatter_combined}) zachycuje simultánně vztah pokrytí \ac{KS} ke čtyřem výstupním proměnným a~umožňuje evaluaci omnibusové hypotézy práce: silnější \ac{KV} instituce jsou systematicky asociovány s~lepšími výsledky napříč všemi sledovanými dimenzemi.
Každý dílčí panel potvrzuje parciální asociaci, ale kombinovaný pohled odhaluje koherentní systémový vzorec: státy s~vysokým pokrytím \ac{KS} (AT, DK, BE, FR) konzistentně vykazují vyšší příjmy v~PPS, nižší pracovní dobu za srovnatelný příjem a~nižší Gini.
ČR se ve všech čtyřech dílčích grafech nachází v~kvadrantu nízkého pokrytí, ačkoli absolutně není na dně EU27.
Polysystémový argument je silnější než jakákoli dílčí korelace: holistické posílení \ac{KV} pokrytí (mechanismus erga omnes, větší dopad \ac{APZ}) by s~vysokou pravděpodobností zlepšilo celou sadu výstupů současně, protože sdílejí společný institucionální determinant.

Kombinovaný rozptylový diagram zachycující zároveň vztah pokrytí kolektivními smlouvami ke čtyřem výstupním proměnným (čistý příjem v~\ac{PPS}, Giniho koeficient, počet odpracovaných hodin, zaměstnanecký podíl vs.~\ac{OSVČ}) umožňuje simultánní evaluaci omnibusové hypotézy práce: silnější \ac{KV} instituce jsou systematicky asociovány s~lepšími výsledky napříč všemi sledovanými dimenzemi.
Každý z~dílčích grafů v~kombinaci potvrzuje parciální asociaci, ale kombinovaný pohled ukazuje, že jde o~koherentní systémový vzorec, nikoliv o~náhodné koincidence: země s~vysokým pokrytím \ac{KS} (AT, DK, BE, FR) konzistentně vykazují vyšší příjmy v~\ac{PPS}, nižší Gini, méně hodin práce za srovnatelný příjem a~nižší podíl \ac{OSVČ}.
ČR se ve všech čtyřech dílčích grafech umísťuje v~kvadrantu „nízké pokrytí / horší výsledky“, ačkoliv v~absolutních hodnotách není na úplném dně EU27.
Tento polysystémový argument je silnější než dílčí korelace: holistické zlepšení institucionálního rámce \ac{KV} (erga omnes + \ac{APZ}) by s~vysokou pravděpodobností zlepšilo celou sadu výstupů současně, protože sdílejí společný institucionální determinant (srov.~obr.~\ref{fig:coverage_income_scatter}, \ref{fig:coverage_gini_scatter}).

Kombinovaný rozptylový diagram zachycující zároveň vztah pokrytí kolektivními smlouvami ke čtyřem výstupním proměnným (čistý příjem v~\ac{PPS}, Giniho koeficient, počet odpracovaných hodin, zaměstnanecký podíl vs.~\ac{OSVČ}) umožňuje simultánní evaluaci omnibusové hypotézy práce: silnější \ac{KV} instituce jsou systematicky asociovány s~lepšími výsledky napříč všemi sledovanými dimenzemi.
Každý z~dílčích grafů v~kombinaci potvrzuje parciální asociaci, ale kombinovaný pohled ukazuje, že jde o~koherentní systémový vzorec, nikoliv o~náhodné koincidence: země s~vysokým pokrytím \ac{KS} (AT, DK, BE, FR) konzistentně vykazují vyšší příjmy v~\ac{PPS}, nižší Gini, méně hodin práce za srovnatelný příjem a~nižší podíl \ac{OSVČ}.
ČR se ve všech čtyřech dílčích grafech umísťuje v~kvadrantu „nízké pokrytí / horší výsledky“, ačkoliv v~absolutních hodnotách není na úplném dně EU27.
Tento polysystémový argument je silnější než dílčí korelace: holistické zlepšení institucionálního rámce \ac{KV} (erga omnes + \ac{APZ}) by s~vysokou pravděpodobností zlepšilo celou sadu výstupů současně, protože sdílejí společný institucionální determinant (srov.~obr.~\ref{fig:coverage_income_scatter}, \ref{fig:coverage_gini_scatter}).

Kombinovaný rozptylový diagram zachycující zároveň vztah pokrytí kolektivními smlouvami ke čtyřem výstupním proměnným (čistý příjem v~\ac{PPS}, Giniho koeficient, počet odpracovaných hodin, zaměstnanecký podíl vs.~\ac{OSVČ}) umožňuje simultánní evaluaci omnibusové hypotézy práce: silnější \ac{KV} instituce jsou systematicky asociovány s~lepšími výsledky napříč všemi sledovanými dimenzemi.
Každý z~dílčích grafů v~kombinaci potvrzuje parciální asociaci, ale kombinovaný pohled ukazuje, že jde o~koherentní systémový vzorec, nikoliv o~náhodné koincidence: země s~vysokým pokrytím \ac{KS} (AT, DK, BE, FR) konzistentně vykazují vyšší příjmy v~\ac{PPS}, nižší Gini, méně hodin práce za srovnatelný příjem a~nižší podíl \ac{OSVČ}.
ČR se ve všech čtyřech dílčích grafech umísťuje v~kvadrantu „nízké pokrytí / horší výsledky“, ačkoliv v~absolutních hodnotách není na úplném dně EU27.
Tento polysystémový argument je silnější než dílčí korelace: holistické zlepšení institucionálního rámce \ac{KV} (erga omnes + \ac{APZ}) by s~vysokou pravděpodobností zlepšilo celou sadu výstupů současně, protože sdílejí společný institucionální determinant (srov.~obr.~\ref{fig:coverage_income_scatter}, \ref{fig:coverage_gini_scatter}).

\input{texparts/python/scatter_combined}




\paragraph{Tabulka korelačních koeficientů}\label{par:corr_table}

Tabulka~\ref{tab:coverage_correlation_table} shrnuje výsledky korelační analýzy za roky 2004--2024.
Každý řádek reportuje: počet párových pozorování~$n$, Pearsonovo~$r$ (rov.~\ref{eq:pearson_r}) na celém panelu, Spearmanovo~$\rho$ (rov.~\ref{eq:spearman_rho}) jako robustní neparametrická alternativa a~$r_{\text{akt.}}$ = Pearsonovo~$r$ pro průřez roku~2024.

Tabulka \ref{tab:coverage_correlation_table} shrnuje výsledky korelační analýzy provedené na párovém panelu zemí OECD/\ac{EU} za dostupná roky od roku 2004.
Zdrojem dat pro pokrytí kolektivního vyjednávání je databáze OECD/AIAS ICTWSS (ukazatel CBC, míra ERB), pro hustotu odborů tatáž databáze (TUD); výstupní proměnné pocházejí z~Eurostatu (lfsa\_ewhun2, nama\_10\_pc, ilc\_di12, earn\_nt\_net).
Každý řádek reportuje: počet párových pozorování $n$ (geo$\times$rok po párování datových sad), Pearsonovo $r$ na celém panelu (míra lineární asociace), Spearmanovo $\rho$ jako robustní neparametrická alternativa rezistentní vůči odlehlým hodnotám, a~$r_{\text{akt.}}$ = Pearsonovo $r$ pro průřezová data z~roku 2024.
Síla korelací je hodnocena dle Evansovy (1996) škály: $|r| < 0{,}20$ velmi slabá; $0{,}20$--$0{,}39$ slabá; $0{,}40$--$0{,}59$ střední; $0{,}60$--$0{,}79$ silná; $\geq 0{,}80$ velmi silná.

Výsledky jsou interpretovány po jednotlivých párech:

\begin{itemize}
  \item \textbf{Pokrytí \ac{KV} vs.~pracovní doba} ($r = -0{,}492$; $\rho = -0{,}384$):
    panelová Pearsonova korelace svědčí o~\emph{střední záporné} asociaci --- vyšší pokrytí \ac{KS} je systematicky spojeno s~kratší průměrnou pracovní dobou.
    Mírný pokles na Spearmanovo $\rho$ (slabá--střední) naznačuje přítomnost několika málo vlivných pozorování posouvajících Pearsonův koeficient, nikoliv výrazný outlier narušující charakter vztahu.
    Průřezová korelace $r_{\text{akt.}} = -0{,}536$ (\emph{střední záporná}) se oproti panelové hodnotě spíše zesílila --- vztah tedy není artefaktem historické variace, ale přetrvává i~v~aktuálním mezinárodním průřezu.

  \item \textbf{Pokrytí \ac{KV} vs.~\ac{HDP}/ob.~(\ac{PPS})} ($r = 0{,}471$; $\rho = 0{,}456$):
    \emph{střední kladná} asociace, vysoce konzistentní mezi Pearsonem a~Spearmanem --- rozdíl pouhých 0,015 indikuje přibližně lineární vztah bez dominantních odlehlých pozorování.
    Průřezová hodnota $r_{\text{akt.}} = 0{,}598$ se přibližuje hranici \emph{silné korelace} (Evans: 0,60), naznačujíce, že v~aktuálním průřezu roku 2024 je vazba mezi institucemi \ac{KV} a~životní úrovní ještě zřetelnější než v~historickém panelu.

  \item \textbf{Pokrytí \ac{KV} vs.~Gini} ($r = -0{,}445$; $\rho = -0{,}286$):
    tento pár vykazuje nejmarkantnější divergenci mezi Pearsonem a~Spearmanem (rozdíl 0,159), která diagnosticky ukazuje na přítomnost outlierů zvyšujících Pearsonův koeficient.
    Nejpravděpodobnějším kandidátem je Francie: s~pokrytím \ac{KS} přes \SI{95}{\percent} (erga omnes) vykazuje relativně vysoké Gini (30--32), protože smluvní komprese mezd je přebita vysokou nerovností kapitálových příjmů a~vrcholných odměn.
    Průřezová hodnota $r_{\text{akt.}} = -0{,}238$ (\emph{slabá záporná}) dále ilustruje, proč nelze z~jednoduchého průřezu soudit o~teorii: v~roce 2024 komprimují Gini napříč EU27 silné transferové systémy a~progresivní zdanění, přičemž \ac{KV} je jen jedním z~více kolineárních institucio\-nálních faktorů.
    Panelový koeficient ($-0{,}445$) sledující \emph{změny} v~čase v~rámci zemí je proto informačně hodnotnější: dokumentuje, že kdy\-koli pokrytí \ac{KS} v~dané zemi kleslo, Gini se tendovalo zvýšit.

  \item \textbf{Pokrytí \ac{KV} vs.~příjem/\ac{HDP}} ($r = 0{,}227$; $\rho = 0{,}230$):
    \emph{slabá kladná} asociace, numericky konzistentní v~obou metodách.
    Relativně nízká hodnota je z~velké části artefaktem konstrukce ukazatele: \ac{HDP}/ob.~je sám o~sobě pozitivně korelován s~\ac{KV} pokrytím ($r = 0{,}471$), takže ve jmenovateli podílu čistý~příjem/\ac{HDP} oba členy rostou společně s~pokrytím a~efekt se vzájemně tlumí.
    Tento výsledek tedy nesvědčí o~tom, že \ac{KV} instituce nezvedají mzdy, ale o~tom, že tento konkrétní relativní ukazatel není citlivým testem dané hypotézy --- absolutní příjmová míra (předchozí pár) zůstává spolehlivějším indikátorem.

  \item \textbf{Hustota odborů vs.~Gini} ($r = -0{,}454$; $\rho = -0{,}455$):
    tento pár je statisticky nejrobustnějším v~celé tabulce: Pearsonovo $r$ a~Spearmanovo $\rho$ jsou prakticky totožné (rozdíl $< 0{,}001$), což svědčí o~přibližně lineárním vztahu bez systematického vlivu odlehlých hodnot.
    \emph{Střední záporná} asociace je stabilní v~celém panelu i~v~aktuálním průřezu ($r_{\text{akt.}} = -0{,}541$).
    Shoda obou koeficientů rovněž zpětně potvrzuje, že divergence u~páru \ac{KV}~vs.~Gini je specifická pro pokrytí (kde Francie tvoří strukturální outlier), nikoliv obecná vlastnost vztahu institucí \ac{KV} k~nerovnosti.
\end{itemize}

Žádná z~výše uvedených korelací nefalzifikuje teorii mzdové komprese prostřednictvím kolektivního vyjednávání.
Klíčovým argumentem je, že jednoduchou bivariátní korelací na heterogenním průřezu 27 zemí lze jen obtížně izolovat marginální příspěvek \ac{KV} pokrytí --- vzdor tomu jsou všechny korelace ve správném předpokládaném směru bez jediné výjimky.
Kauzální identifikaci přináší panelová analýza s~fixními efekty a~přirozené experimenty: metaanalytická literatura (např.~Jaumotte \& Osorio Buitron 2015; Card et al.~2004) konzistentně dokládá statisticky i~věcně významný efekt \ac{KV} institucí na mzdovou kompresi.
Korelace v~tabulce tedy plní funkci deskriptivního podkladu motivujícího institucionální případ\-ovou studii ČR --- nikoliv jako izolovaný kauzální důkaz, ale jako konzistentní mezinárodní vzorec, v~němž ČR se svým pokrytím $\approx 32\,\%$ systematicky vykazuje horšíu výsledky než srovnatelné ekonomiky s~vyšším pokrytím \ac{KS}.

Tabulka \ref{tab:coverage_correlation_table} shrnuje výsledky korelační analýzy provedené na párovém panelu zemí OECD/\ac{EU} za dostupná roky od roku 2004.
Zdrojem dat pro pokrytí kolektivního vyjednávání je databáze OECD/AIAS ICTWSS (ukazatel CBC, míra ERB), pro hustotu odborů tatáž databáze (TUD); výstupní proměnné pocházejí z~Eurostatu (lfsa\_ewhun2, nama\_10\_pc, ilc\_di12, earn\_nt\_net).
Každý řádek reportuje: počet párových pozorování $n$ (geo$\times$rok po párování datových sad), Pearsonovo $r$ na celém panelu (míra lineární asociace), Spearmanovo $\rho$ jako robustní neparametrická alternativa rezistentní vůči odlehlým hodnotám, a~$r_{\text{akt.}}$ = Pearsonovo $r$ pro průřezová data z~roku 2024.
Síla korelací je hodnocena dle Evansovy (1996) škály: $|r| < 0{,}20$ velmi slabá; $0{,}20$--$0{,}39$ slabá; $0{,}40$--$0{,}59$ střední; $0{,}60$--$0{,}79$ silná; $\geq 0{,}80$ velmi silná.

Výsledky jsou interpretovány po jednotlivých párech:

\begin{itemize}
  \item \textbf{Pokrytí \ac{KV} vs.~pracovní doba} ($r = -0{,}492$; $\rho = -0{,}384$):
    panelová Pearsonova korelace svědčí o~\emph{střední záporné} asociaci --- vyšší pokrytí \ac{KS} je systematicky spojeno s~kratší průměrnou pracovní dobou.
    Mírný pokles na Spearmanovo $\rho$ (slabá--střední) naznačuje přítomnost několika málo vlivných pozorování posouvajících Pearsonův koeficient, nikoliv výrazný outlier narušující charakter vztahu.
    Průřezová korelace $r_{\text{akt.}} = -0{,}536$ (\emph{střední záporná}) se oproti panelové hodnotě spíše zesílila --- vztah tedy není artefaktem historické variace, ale přetrvává i~v~aktuálním mezinárodním průřezu.

  \item \textbf{Pokrytí \ac{KV} vs.~\ac{HDP}/ob.~(\ac{PPS})} ($r = 0{,}471$; $\rho = 0{,}456$):
    \emph{střední kladná} asociace, vysoce konzistentní mezi Pearsonem a~Spearmanem --- rozdíl pouhých 0,015 indikuje přibližně lineární vztah bez dominantních odlehlých pozorování.
    Průřezová hodnota $r_{\text{akt.}} = 0{,}598$ se přibližuje hranici \emph{silné korelace} (Evans: 0,60), naznačujíce, že v~aktuálním průřezu roku 2024 je vazba mezi institucemi \ac{KV} a~životní úrovní ještě zřetelnější než v~historickém panelu.

  \item \textbf{Pokrytí \ac{KV} vs.~Gini} ($r = -0{,}445$; $\rho = -0{,}286$):
    tento pár vykazuje nejmarkantnější divergenci mezi Pearsonem a~Spearmanem (rozdíl 0,159), která diagnosticky ukazuje na přítomnost outlierů zvyšujících Pearsonův koeficient.
    Nejpravděpodobnějším kandidátem je Francie: s~pokrytím \ac{KS} přes \SI{95}{\percent} (erga omnes) vykazuje relativně vysoké Gini (30--32), protože smluvní komprese mezd je přebita vysokou nerovností kapitálových příjmů a~vrcholných odměn.
    Průřezová hodnota $r_{\text{akt.}} = -0{,}238$ (\emph{slabá záporná}) dále ilustruje, proč nelze z~jednoduchého průřezu soudit o~teorii: v~roce 2024 komprimují Gini napříč EU27 silné transferové systémy a~progresivní zdanění, přičemž \ac{KV} je jen jedním z~více kolineárních institucio\-nálních faktorů.
    Panelový koeficient ($-0{,}445$) sledující \emph{změny} v~čase v~rámci zemí je proto informačně hodnotnější: dokumentuje, že kdy\-koli pokrytí \ac{KS} v~dané zemi kleslo, Gini se tendovalo zvýšit.

  \item \textbf{Pokrytí \ac{KV} vs.~příjem/\ac{HDP}} ($r = 0{,}227$; $\rho = 0{,}230$):
    \emph{slabá kladná} asociace, numericky konzistentní v~obou metodách.
    Relativně nízká hodnota je z~velké části artefaktem konstrukce ukazatele: \ac{HDP}/ob.~je sám o~sobě pozitivně korelován s~\ac{KV} pokrytím ($r = 0{,}471$), takže ve jmenovateli podílu čistý~příjem/\ac{HDP} oba členy rostou společně s~pokrytím a~efekt se vzájemně tlumí.
    Tento výsledek tedy nesvědčí o~tom, že \ac{KV} instituce nezvedají mzdy, ale o~tom, že tento konkrétní relativní ukazatel není citlivým testem dané hypotézy --- absolutní příjmová míra (předchozí pár) zůstává spolehlivějším indikátorem.

  \item \textbf{Hustota odborů vs.~Gini} ($r = -0{,}454$; $\rho = -0{,}455$):
    tento pár je statisticky nejrobustnějším v~celé tabulce: Pearsonovo $r$ a~Spearmanovo $\rho$ jsou prakticky totožné (rozdíl $< 0{,}001$), což svědčí o~přibližně lineárním vztahu bez systematického vlivu odlehlých hodnot.
    \emph{Střední záporná} asociace je stabilní v~celém panelu i~v~aktuálním průřezu ($r_{\text{akt.}} = -0{,}541$).
    Shoda obou koeficientů rovněž zpětně potvrzuje, že divergence u~páru \ac{KV}~vs.~Gini je specifická pro pokrytí (kde Francie tvoří strukturální outlier), nikoliv obecná vlastnost vztahu institucí \ac{KV} k~nerovnosti.
\end{itemize}

Žádná z~výše uvedených korelací nefalzifikuje teorii mzdové komprese prostřednictvím kolektivního vyjednávání.
Klíčovým argumentem je, že jednoduchou bivariátní korelací na heterogenním průřezu 27 zemí lze jen obtížně izolovat marginální příspěvek \ac{KV} pokrytí --- vzdor tomu jsou všechny korelace ve správném předpokládaném směru bez jediné výjimky.
Kauzální identifikaci přináší panelová analýza s~fixními efekty a~přirozené experimenty: metaanalytická literatura (např.~Jaumotte \& Osorio Buitron 2015; Card et al.~2004) konzistentně dokládá statisticky i~věcně významný efekt \ac{KV} institucí na mzdovou kompresi.
Korelace v~tabulce tedy plní funkci deskriptivního podkladu motivujícího institucionální případ\-ovou studii ČR --- nikoliv jako izolovaný kauzální důkaz, ale jako konzistentní mezinárodní vzorec, v~němž ČR se svým pokrytím $\approx 32\,\%$ systematicky vykazuje horšíu výsledky než srovnatelné ekonomiky s~vyšším pokrytím \ac{KS}.

Tabulka \ref{tab:coverage_correlation_table} shrnuje výsledky korelační analýzy provedené na párovém panelu zemí OECD/\ac{EU} za dostupná roky od roku 2004.
Zdrojem dat pro pokrytí kolektivního vyjednávání je databáze OECD/AIAS ICTWSS (ukazatel CBC, míra ERB), pro hustotu odborů tatáž databáze (TUD); výstupní proměnné pocházejí z~Eurostatu (lfsa\_ewhun2, nama\_10\_pc, ilc\_di12, earn\_nt\_net).
Každý řádek reportuje: počet párových pozorování $n$ (geo$\times$rok po párování datových sad), Pearsonovo $r$ na celém panelu (míra lineární asociace), Spearmanovo $\rho$ jako robustní neparametrická alternativa rezistentní vůči odlehlým hodnotám, a~$r_{\text{akt.}}$ = Pearsonovo $r$ pro průřezová data z~roku 2024.
Síla korelací je hodnocena dle Evansovy (1996) škály: $|r| < 0{,}20$ velmi slabá; $0{,}20$--$0{,}39$ slabá; $0{,}40$--$0{,}59$ střední; $0{,}60$--$0{,}79$ silná; $\geq 0{,}80$ velmi silná.

Výsledky jsou interpretovány po jednotlivých párech:

\begin{itemize}
  \item \textbf{Pokrytí \ac{KV} vs.~pracovní doba} ($r = -0{,}492$; $\rho = -0{,}384$):
    panelová Pearsonova korelace svědčí o~\emph{střední záporné} asociaci --- vyšší pokrytí \ac{KS} je systematicky spojeno s~kratší průměrnou pracovní dobou.
    Mírný pokles na Spearmanovo $\rho$ (slabá--střední) naznačuje přítomnost několika málo vlivných pozorování posouvajících Pearsonův koeficient, nikoliv výrazný outlier narušující charakter vztahu.
    Průřezová korelace $r_{\text{akt.}} = -0{,}536$ (\emph{střední záporná}) se oproti panelové hodnotě spíše zesílila --- vztah tedy není artefaktem historické variace, ale přetrvává i~v~aktuálním mezinárodním průřezu.

  \item \textbf{Pokrytí \ac{KV} vs.~\ac{HDP}/ob.~(\ac{PPS})} ($r = 0{,}471$; $\rho = 0{,}456$):
    \emph{střední kladná} asociace, vysoce konzistentní mezi Pearsonem a~Spearmanem --- rozdíl pouhých 0,015 indikuje přibližně lineární vztah bez dominantních odlehlých pozorování.
    Průřezová hodnota $r_{\text{akt.}} = 0{,}598$ se přibližuje hranici \emph{silné korelace} (Evans: 0,60), naznačujíce, že v~aktuálním průřezu roku 2024 je vazba mezi institucemi \ac{KV} a~životní úrovní ještě zřetelnější než v~historickém panelu.

  \item \textbf{Pokrytí \ac{KV} vs.~Gini} ($r = -0{,}445$; $\rho = -0{,}286$):
    tento pár vykazuje nejmarkantnější divergenci mezi Pearsonem a~Spearmanem (rozdíl 0,159), která diagnosticky ukazuje na přítomnost outlierů zvyšujících Pearsonův koeficient.
    Nejpravděpodobnějším kandidátem je Francie: s~pokrytím \ac{KS} přes \SI{95}{\percent} (erga omnes) vykazuje relativně vysoké Gini (30--32), protože smluvní komprese mezd je přebita vysokou nerovností kapitálových příjmů a~vrcholných odměn.
    Průřezová hodnota $r_{\text{akt.}} = -0{,}238$ (\emph{slabá záporná}) dále ilustruje, proč nelze z~jednoduchého průřezu soudit o~teorii: v~roce 2024 komprimují Gini napříč EU27 silné transferové systémy a~progresivní zdanění, přičemž \ac{KV} je jen jedním z~více kolineárních institucio\-nálních faktorů.
    Panelový koeficient ($-0{,}445$) sledující \emph{změny} v~čase v~rámci zemí je proto informačně hodnotnější: dokumentuje, že kdy\-koli pokrytí \ac{KS} v~dané zemi kleslo, Gini se tendovalo zvýšit.

  \item \textbf{Pokrytí \ac{KV} vs.~příjem/\ac{HDP}} ($r = 0{,}227$; $\rho = 0{,}230$):
    \emph{slabá kladná} asociace, numericky konzistentní v~obou metodách.
    Relativně nízká hodnota je z~velké části artefaktem konstrukce ukazatele: \ac{HDP}/ob.~je sám o~sobě pozitivně korelován s~\ac{KV} pokrytím ($r = 0{,}471$), takže ve jmenovateli podílu čistý~příjem/\ac{HDP} oba členy rostou společně s~pokrytím a~efekt se vzájemně tlumí.
    Tento výsledek tedy nesvědčí o~tom, že \ac{KV} instituce nezvedají mzdy, ale o~tom, že tento konkrétní relativní ukazatel není citlivým testem dané hypotézy --- absolutní příjmová míra (předchozí pár) zůstává spolehlivějším indikátorem.

  \item \textbf{Hustota odborů vs.~Gini} ($r = -0{,}454$; $\rho = -0{,}455$):
    tento pár je statisticky nejrobustnějším v~celé tabulce: Pearsonovo $r$ a~Spearmanovo $\rho$ jsou prakticky totožné (rozdíl $< 0{,}001$), což svědčí o~přibližně lineárním vztahu bez systematického vlivu odlehlých hodnot.
    \emph{Střední záporná} asociace je stabilní v~celém panelu i~v~aktuálním průřezu ($r_{\text{akt.}} = -0{,}541$).
    Shoda obou koeficientů rovněž zpětně potvrzuje, že divergence u~páru \ac{KV}~vs.~Gini je specifická pro pokrytí (kde Francie tvoří strukturální outlier), nikoliv obecná vlastnost vztahu institucí \ac{KV} k~nerovnosti.
\end{itemize}

Žádná z~výše uvedených korelací nefalzifikuje teorii mzdové komprese prostřednictvím kolektivního vyjednávání.
Klíčovým argumentem je, že jednoduchou bivariátní korelací na heterogenním průřezu 27 zemí lze jen obtížně izolovat marginální příspěvek \ac{KV} pokrytí --- vzdor tomu jsou všechny korelace ve správném předpokládaném směru bez jediné výjimky.
Kauzální identifikaci přináší panelová analýza s~fixními efekty a~přirozené experimenty: metaanalytická literatura (např.~Jaumotte \& Osorio Buitron 2015; Card et al.~2004) konzistentně dokládá statisticky i~věcně významný efekt \ac{KV} institucí na mzdovou kompresi.
Korelace v~tabulce tedy plní funkci deskriptivního podkladu motivujícího institucionální případ\-ovou studii ČR --- nikoliv jako izolovaný kauzální důkaz, ale jako konzistentní mezinárodní vzorec, v~němž ČR se svým pokrytím $\approx 32\,\%$ systematicky vykazuje horšíu výsledky než srovnatelné ekonomiky s~vyšším pokrytím \ac{KS}.

\input{texparts/python/coverage_correlation_table}




Výsledky jsou interpretovány po párech:
\begin{description}
  \item[Pokrytí \ac{KV} vs.\ pracovní doba]
    ($r = -0{,}483$; $\rho = -0{,}525$; $R^2 = 0{,}23$; $n = 256$):
    střední záporná asociace robustně potvrzená oběma metodami.
    Spearmanovo $\rho$ je v~absolutní hodnotě vyšší než Pearsonovo~$r$, což svědčí o~monotónním charakteru vztahu bez dominantních outlierů posouvajících lineární složku.
    Průřezová korelace $r_{\text{akt.}} = -0{,}324$ je slabší než panelová hodnota, pravděpodobně vlivem strukturní konvergence pracovní doby v~rámci EU v~posledním desetiletí (zkracování doby v~CEE, stagnace na severu).

  \item[Pokrytí \ac{KV} vs.\ HDP/ob.\ (PPS)]
    ($r = 0{,}562$; $\rho = 0{,}450$; $R^2 = 0{,}32$; $n = 299$):
    střední až silná kladná asociace; Pearsonovo~$r$ je o~0,112 vyšší než Spearmanovo~$\rho$, což naznačuje citlivost panelového odhadu na severské země s~extrémně vysokým pokrytím i~HDP.
    Vyřazení Irska posílilo panelový signál (Irsko kombinovalo nízké pokrytí s~extrémním HDP/ob., což dříve tlumilo kladnou korelaci); průřezová hodnota $r_{\text{akt.}} = 0{,}448$ zůstává střední.

  \item[Pokrytí \ac{KV} vs.\ Gini]
    ($r = -0{,}303$; $\rho = -0{,}238$; $R^2 = 0{,}09$; $n = 163$):
    slabá záporná asociace; divergence $|r - \rho| = 0{,}065$ je výrazně menší než při zahrnutí Irska (tehdy 0,159), čímž je potvrzeno, že Irsko bylo strukturálním outlierem narušujícím tuto korelaci.
    Průřezová hodnota $r_{\text{akt.}} = +0{,}481$ je kladná, neboť v~průřezu roku~2024 mají státy s~erga omnes pokrytím (FR, IT, BE, AT --- všechny $\geq 90\,\%$) vyšší disponibilní Gini díky vysokým kapitálovým příjmům a~odměnám vrcholového managementu, nikoliv absenci mzdové komprese.
    Panelový koeficient sledující \emph{změny v~čase v~rámci zemí} je informačně hodnotnější: dokumentuje, že kdykoli pokrytí \ac{KS} v~dané zemi kleslo, Gini se tendovalo zvýšit.
    Příčinný mechanismus KV\,$\to$\,Gini operuje zejména přes primární mzdovou distribuci, nikoliv přes disponibilní příjmy ovlivněné fiskální redistribucí.

  \item[Pokrytí \ac{KV} vs.\ příjem/HDP]
    ($r = 0{,}210$; $\rho = 0{,}177$; $R^2 = 0{,}04$; $n = 286$):
    slabá kladná asociace, konzistentní v~obou metodách.
    Nízká hodnota je z~velké části artefaktem konstrukce ukazatele: HDP/ob.~je sám středně silně korelováno s~pokrytím ($r = 0{,}562$), takže ve jmenovateli podílu oba členy rostou společně s~pokrytím a~efekt se vzájemně tlumí.
    Absolutní příjmová míra (předchozí pár) je proto citlivějším testem hypotézy o~mzdovém efektu \ac{KV}.

  \item[Hustota odborů vs.\ Gini]
    ($r = -0{,}458$; $\rho = -0{,}462$):
    statisticky nejrobustnější pár v~celé tabulce: Pearsonovo~$r$ a~Spearmanovo~$\rho$ jsou prakticky totožné (rozdíl 0,004), což svědčí o~přibližně lineárním vztahu bez systematického vlivu odlehlých hodnot.
    Průřezová korelace $r_{\text{akt.}} = -0{,}627$ dosahuje \emph{silné záporné} asociace a~je nejsilnějším institucionálním prediktorem nízkého Gini ze všech testovaných párů.
    Shoda panelových koeficientů a~jejich posílení v~průřezu zpětně potvrzuje, že kladná průřezová hodnota u~páru pokrytí~vs.~Gini je specifická pro pokrytí (kde státy jako Francie deformují průřezový vztah), nikoliv obecná vlastnost vztahu \ac{KV} institucí k~nerovnosti.
\end{description}

Žádná z~panelových korelací nefalzifikuje teorii mzdové komprese prostřednictvím \ac{KV} institucí.
Jednoduchou bivariátní korelací na heterogenním průřezu 24~zemí nelze izolovat marginální příspěvek pokrytí \ac{KS}; přesto jsou všechny panelové korelace ve správném předpokládaném směru bez výjimky.
Kauzální identifikaci přináší panelová analýza s~fixními efekty: metaanalytická literatura (Jaumotte \& Osorio Buitron, 2015; Card a~kol., 2004) konzistentně dokládá věcně i~statisticky významný efekt \ac{KV} institucí na mzdovou kompresi.
Korelace v~tabulce tak plní funkci deskriptivního podkladu motivujícího institucionální případovou studii ČR --- jako konzistentní mezinárodní vzorec, v~němž ČR se svým pokrytím $\approx 31\,\%$ systematicky vykazuje horší výsledky než srovnatelné ekonomiky s~vyšším pokrytím \ac{KS}.

\paragraph{Pokrytí \ac{KS} a pracovní doba}\label{par:hours_scatter}

Bodový graf (obr.~\ref{fig:coverage_hours_scatter}) ukazuje zápornou korelaci mezi pokrytím \ac{KS} a~průměrnou týdenní pracovní dobou ($r = -0{,}483$; $\hat{\beta} = -0{,}042$~h/týden na pp pokrytí dle rov.~\ref{eq:ols}): pracovníci v~zemích s~vyšším pokrytím pracují v~průměru méně hodin při srovnatelném nebo vyšším reálném příjmu.
Efekt je konzistentní s~teorií, podle níž silné \ac{KS} prosazují pracovní dobu jako plnohodnotnou vyjednávací proměnnou --- zkrácené úvazky, limit přesčasů, minimální odpočinkové doby nad zákonné standardy.
V~ČR, kde \ac{KS} zřídka obsahují ustanovení o~pracovní době nad zákonné minimum, pracovníci ve srovnatelných odvětvích odpracují více hodin než v~AT nebo DK, přičemž hodinová mzdová sazba je přesto nižší.
Prodloužená pracovní doba v~ČR (nadprůměrný podíl přesčasů, nejnižší podíl částečných úvazků v~EU) nesignalizuje vyšší produktivitu, ale neschopnost \ac{KS} prosadit kompenzaci výkonu formou kratší doby nebo vyšší sazby za přesčas.
Reforma \ac{KV} by proto měla výslovně zahrnout pracovní dobu jako standardní vyjednávací agendu: \ac{KS}, která řeší jen základní mzdu bez úpravy doby, nevyjednává celkovou hodnotu pracovního vztahu.

Bodový graf závislosti průměrného počtu odpracovaných hodin za rok na pokrytí \ac{KS} ukazuje zápornou korelaci: pracovníci v~zemích s~vyšším \ac{KS} pokrytím pracují v~průměru méně hodin ročně při srovnatelném nebo vyšším reálném příjmu --- efekt, který je konzistentní s~teorií, podle níž silné \ac{KS} prosazují i~pracovní dobu jako vyjednávací proměnnou (zkrácené pracovní úvazky, limit přesčasů).
V~ČR, kde \ac{KS} zřídka obsahují robustní ustanovení o~pracovní době nad zákonné minimum, pracovníci ve srovnatelných odvětvích odpracují více hodin než v~AT nebo DK, přičemž hodinová mzdová sazba je přesto nižší --- kombinace více hodin za méně peněz je symptomem absence koordinovaného vyjednávání.
Prodloužená pracovní doba v~ČR (nadprůměrný podíl přesčasů, nejnižší podíl částečných úvazků v~\ac{EU}) nesignalizuje vyšší pracovní nasazení, ale neschopnost kolektivní smlouvy prosadit kompenzaci produktivity formou kratší doby nebo vyšší sazby za přesčas.
Reforma \ac{KV} by proto měla výslovně zahrnout pracovní dobu jako standardní vyjednávací agendu: \ac{KS}, která řeší jen základní mzdu bez doby, nevyjednává celkovou hodnotu pracovního vztahu.

Bodový graf závislosti průměrného počtu odpracovaných hodin za rok na pokrytí \ac{KS} ukazuje zápornou korelaci: pracovníci v~zemích s~vyšším \ac{KS} pokrytím pracují v~průměru méně hodin ročně při srovnatelném nebo vyšším reálném příjmu --- efekt, který je konzistentní s~teorií, podle níž silné \ac{KS} prosazují i~pracovní dobu jako vyjednávací proměnnou (zkrácené pracovní úvazky, limit přesčasů).
V~ČR, kde \ac{KS} zřídka obsahují robustní ustanovení o~pracovní době nad zákonné minimum, pracovníci ve srovnatelných odvětvích odpracují více hodin než v~AT nebo DK, přičemž hodinová mzdová sazba je přesto nižší --- kombinace více hodin za méně peněz je symptomem absence koordinovaného vyjednávání.
Prodloužená pracovní doba v~ČR (nadprůměrný podíl přesčasů, nejnižší podíl částečných úvazků v~\ac{EU}) nesignalizuje vyšší pracovní nasazení, ale neschopnost kolektivní smlouvy prosadit kompenzaci produktivity formou kratší doby nebo vyšší sazby za přesčas.
Reforma \ac{KV} by proto měla výslovně zahrnout pracovní dobu jako standardní vyjednávací agendu: \ac{KS}, která řeší jen základní mzdu bez doby, nevyjednává celkovou hodnotu pracovního vztahu.

Bodový graf závislosti průměrného počtu odpracovaných hodin za rok na pokrytí \ac{KS} ukazuje zápornou korelaci: pracovníci v~zemích s~vyšším \ac{KS} pokrytím pracují v~průměru méně hodin ročně při srovnatelném nebo vyšším reálném příjmu --- efekt, který je konzistentní s~teorií, podle níž silné \ac{KS} prosazují i~pracovní dobu jako vyjednávací proměnnou (zkrácené pracovní úvazky, limit přesčasů).
V~ČR, kde \ac{KS} zřídka obsahují robustní ustanovení o~pracovní době nad zákonné minimum, pracovníci ve srovnatelných odvětvích odpracují více hodin než v~AT nebo DK, přičemž hodinová mzdová sazba je přesto nižší --- kombinace více hodin za méně peněz je symptomem absence koordinovaného vyjednávání.
Prodloužená pracovní doba v~ČR (nadprůměrný podíl přesčasů, nejnižší podíl částečných úvazků v~\ac{EU}) nesignalizuje vyšší pracovní nasazení, ale neschopnost kolektivní smlouvy prosadit kompenzaci produktivity formou kratší doby nebo vyšší sazby za přesčas.
Reforma \ac{KV} by proto měla výslovně zahrnout pracovní dobu jako standardní vyjednávací agendu: \ac{KS}, která řeší jen základní mzdu bez doby, nevyjednává celkovou hodnotu pracovního vztahu.

\input{texparts/python/coverage_hours_scatter}




\paragraph{Pokrytí \ac{KS} a HDP na obyvatele v~PPS}\label{par:income_scatter}

Bodový graf (obr.~\ref{fig:coverage_income_scatter}) dokumentuje kladnou korelaci pokrytí \ac{KS} s~HDP na obyvatele v~PPS ($r = 0{,}562$; $\hat{\beta} = 0{,}476$~indexových bodů na pp pokrytí; $R^2 = 0{,}316$).
Kauzální interpretaci podporuje teorie vyjednávání: \ac{KS} s~plošnou platností (erga omnes) komprimují mzdové rozdíly a~zvyšují spodní část mzdové distribuce, čímž přispívají k~vyšší průměrné životní úrovni v~celé ekonomice.
ČR se nachází nalevo od středu grafu (pokrytí $\approx 31\,\%$, HDP/ob.\ index~92): hypotetické přiblížení k~cílové hodnotě 80~\% by dle regresní přímky odpovídalo nárůstu indexu HDP/ob.\ o~přibližně 23~bodů, čímž by ČR překročila průměr EU27.
Tento obrázek slouží jako ústřední vizualizace teze práce: posílení \ac{KV} pokrytí není jen otázkou pracovního práva, ale klíčovým makroekonomickým nástrojem pro zvýšení životní úrovně bez nutnosti přímé státní redistribuce.

Bodový graf závislosti čistého příjmu domácností v~\ac{PPS} na pokrytí \ac{KS} je stěžejním grafem této práce: dokumentuje kladnou korelaci, podle níž státy s~vyšším pokrytím \ac{KS} vykazují vyšší reálné čisté příjmy --- a~to i~po kontrole na úroveň \ac{HDP} na obyvatele.
Kauzální interpretace tohoto vztahu je podpořena teorií vyjednávání: \ac{KS} s~plošnou platností komprimují mzdové rozdíly a~zvyšují spodní část mzdové distribuce, čímž přispívají k~vyšší průměrné a~mediánové úrovni příjmu v~celé ekonomice.
ČR se umísťuje nalevo od středu grafu s~relativně nízkým příjmem v~\ac{PPS} a~nízkým \ac{KS} pokrytím (cca \SI{35}{\percent}): hypotetické přiblížení k~pokrytí \SIrange{70}{80}{\percent} (AT úroveň) by dle korelační linie odpovídalo nárůstu \ac{PPS} příjmu o~\SIrange{15}{20}{\percent} --- efekt srovnatelný s~výstupy z~přirozených experimentů v~literatuře.
Tento obrázek může sloužit jako ústřední vizualizace teze: posílení \ac{KV} pokrytí není jen otázkou pracovního práva, ale klíčovým makroekonomickým nástrojem pro zvýšení životní úrovně bez nutnosti přímé státní redistribuce.

Bodový graf závislosti čistého příjmu domácností v~\ac{PPS} na pokrytí \ac{KS} je stěžejním grafem této práce: dokumentuje kladnou korelaci, podle níž státy s~vyšším pokrytím \ac{KS} vykazují vyšší reálné čisté příjmy --- a~to i~po kontrole na úroveň \ac{HDP} na obyvatele.
Kauzální interpretace tohoto vztahu je podpořena teorií vyjednávání: \ac{KS} s~plošnou platností komprimují mzdové rozdíly a~zvyšují spodní část mzdové distribuce, čímž přispívají k~vyšší průměrné a~mediánové úrovni příjmu v~celé ekonomice.
ČR se umísťuje nalevo od středu grafu s~relativně nízkým příjmem v~\ac{PPS} a~nízkým \ac{KS} pokrytím (cca \SI{35}{\percent}): hypotetické přiblížení k~pokrytí \SIrange{70}{80}{\percent} (AT úroveň) by dle korelační linie odpovídalo nárůstu \ac{PPS} příjmu o~\SIrange{15}{20}{\percent} --- efekt srovnatelný s~výstupy z~přirozených experimentů v~literatuře.
Tento obrázek může sloužit jako ústřední vizualizace teze: posílení \ac{KV} pokrytí není jen otázkou pracovního práva, ale klíčovým makroekonomickým nástrojem pro zvýšení životní úrovně bez nutnosti přímé státní redistribuce.

Bodový graf závislosti čistého příjmu domácností v~\ac{PPS} na pokrytí \ac{KS} je stěžejním grafem této práce: dokumentuje kladnou korelaci, podle níž státy s~vyšším pokrytím \ac{KS} vykazují vyšší reálné čisté příjmy --- a~to i~po kontrole na úroveň \ac{HDP} na obyvatele.
Kauzální interpretace tohoto vztahu je podpořena teorií vyjednávání: \ac{KS} s~plošnou platností komprimují mzdové rozdíly a~zvyšují spodní část mzdové distribuce, čímž přispívají k~vyšší průměrné a~mediánové úrovni příjmu v~celé ekonomice.
ČR se umísťuje nalevo od středu grafu s~relativně nízkým příjmem v~\ac{PPS} a~nízkým \ac{KS} pokrytím (cca \SI{35}{\percent}): hypotetické přiblížení k~pokrytí \SIrange{70}{80}{\percent} (AT úroveň) by dle korelační linie odpovídalo nárůstu \ac{PPS} příjmu o~\SIrange{15}{20}{\percent} --- efekt srovnatelný s~výstupy z~přirozených experimentů v~literatuře.
Tento obrázek může sloužit jako ústřední vizualizace teze: posílení \ac{KV} pokrytí není jen otázkou pracovního práva, ale klíčovým makroekonomickým nástrojem pro zvýšení životní úrovně bez nutnosti přímé státní redistribuce.

\input{texparts/python/coverage_income_scatter}




\paragraph{Pokrytí \ac{KS} a Giniho koeficient}\label{par:gini_scatter}

Bodový graf (obr.~\ref{fig:coverage_gini_scatter}) potvrzuje zápornou panelovou korelaci pokrytí \ac{KS} s~příjmovou nerovností ($r = -0{,}303$; $\hat{\beta} = -0{,}037$~bodu Gini na pp pokrytí).
Mechanismus je přímočarý: odvětvové \ac{KS} závazné pro celý sektor stlačují spodní část mzdové distribuce, zvyšují mzdy pracovníků ve spodním kvartilem a~tím snižují Giniho koeficient.
ČR s~pokrytím $\approx 31\,\%$ a~Gini~24 se nachází v~prostoru grafu, kde by hypotetické zvýšení pokrytí na~80~\% odpovídalo poklesu Gini o~přibližně 1,8~bodu --- změna, která by ČR přiblížila vzoru skandinávských zemí, a~to čistě tržní redistribucí bez zvyšování daní.
Redistributivní efekt \ac{KV} institucí je empiricky prokazatelný i~bez zvyšování přímých daní nebo transferů: jde o~distribuci prostřednictvím primárních mezd.

Bodový graf závislosti Giniho koeficientu příjmové nerovnosti na pokrytí \ac{KS} potvrzuje zápornou korelaci: vyšší \ac{KS} pokrytí je asociováno s~nižší příjmovou nerovností --- vztah, který je v~empirické literatuře o~institucích trhu práce považován za jeden z~nejrobustnějších.
Mechanismus je přímočarý: odvětvové \ac{KS} s~tarify závaznými pro celý sektor (erga omnes) stlačují spodní část mzdové distribuce, zvyšují mezd zaměstnanců ve spodním kvartilem, a~tím snižují roztažení distribuce měřené Ginim.
ČR s~pokrytím cca \SI{35}{\percent} a~aktuálním Gini cca 27 se nachází v~prostoru grafu, kde by hypotetické zvýšení pokrytí na \SI{70}{\percent} (středoevropský cíl) odpovídalo poklesu Gini o~2--3 body --- změna, která by CZ přiblížila finskému nebo belgickému vzoru.
Redistributivní efekt \ac{KV} institucí je tedy empiricky prokazatelný i~bez zvyšování přímých daní nebo transferů: jde o~tržní redistribuci prostřednictvím vyjednávání, která je z~hlediska politické ekonomie snáze prosaditelná než fiskální redistribuce.

Bodový graf závislosti Giniho koeficientu příjmové nerovnosti na pokrytí \ac{KS} potvrzuje zápornou korelaci: vyšší \ac{KS} pokrytí je asociováno s~nižší příjmovou nerovností --- vztah, který je v~empirické literatuře o~institucích trhu práce považován za jeden z~nejrobustnějších.
Mechanismus je přímočarý: odvětvové \ac{KS} s~tarify závaznými pro celý sektor (erga omnes) stlačují spodní část mzdové distribuce, zvyšují mezd zaměstnanců ve spodním kvartilem, a~tím snižují roztažení distribuce měřené Ginim.
ČR s~pokrytím cca \SI{35}{\percent} a~aktuálním Gini cca 27 se nachází v~prostoru grafu, kde by hypotetické zvýšení pokrytí na \SI{70}{\percent} (středoevropský cíl) odpovídalo poklesu Gini o~2--3 body --- změna, která by CZ přiblížila finskému nebo belgickému vzoru.
Redistributivní efekt \ac{KV} institucí je tedy empiricky prokazatelný i~bez zvyšování přímých daní nebo transferů: jde o~tržní redistribuci prostřednictvím vyjednávání, která je z~hlediska politické ekonomie snáze prosaditelná než fiskální redistribuce.

Bodový graf závislosti Giniho koeficientu příjmové nerovnosti na pokrytí \ac{KS} potvrzuje zápornou korelaci: vyšší \ac{KS} pokrytí je asociováno s~nižší příjmovou nerovností --- vztah, který je v~empirické literatuře o~institucích trhu práce považován za jeden z~nejrobustnějších.
Mechanismus je přímočarý: odvětvové \ac{KS} s~tarify závaznými pro celý sektor (erga omnes) stlačují spodní část mzdové distribuce, zvyšují mezd zaměstnanců ve spodním kvartilem, a~tím snižují roztažení distribuce měřené Ginim.
ČR s~pokrytím cca \SI{35}{\percent} a~aktuálním Gini cca 27 se nachází v~prostoru grafu, kde by hypotetické zvýšení pokrytí na \SI{70}{\percent} (středoevropský cíl) odpovídalo poklesu Gini o~2--3 body --- změna, která by CZ přiblížila finskému nebo belgickému vzoru.
Redistributivní efekt \ac{KV} institucí je tedy empiricky prokazatelný i~bez zvyšování přímých daní nebo transferů: jde o~tržní redistribuci prostřednictvím vyjednávání, která je z~hlediska politické ekonomie snáze prosaditelná než fiskální redistribuce.

\input{texparts/python/coverage_gini_scatter}




\paragraph{Pokrytí \ac{KS} a relativní příjem}\label{par:ratio_scatter}

Bodový graf (obr.~\ref{fig:coverage_ratio_scatter}) zachycuje podíl čistého příjmu v~PPS na HDP/ob.~v~PPS jako funkci pokrytí \ac{KS} ($r = 0{,}210$; $R^2 = 0{,}04$).
Slabá asociace je, jak bylo vysvětleno výše, z~velké části artefaktem endogenity jmenovatele: HDP/ob.~je sám středně korelováno s~pokrytím, čímž se efekt ve jmenovateli tlumí.
Klíčovým přínosem tohoto grafu proto není sklon regresní přímky, ale vzájemná poloha zemí: státy s~nízkým podílem čistého příjmu na HDP \emph{a~zároveň} nízkým HDP (ČR, SK, HU) signalizují dvojí selhání trhu práce --- nízká produktivita doprovázená nízkým podílem práce na distribuci vytvořené hodnoty.
Kombinace více hodin~$\times$~nižší mzda~$\times$~nižší podíl na HDP je symptomem ekonomiky postavené na modelu levné práce bez koordinace mezd.
Minimální mzda a~\ac{KS} pokrytí jsou v~tomto kontextu substituty: v~zemích s~vysokým erga omnes pokrytím (DK, BE) zákonná minimální mzda prakticky neplní funkci, protože smluvní minima jsou nastavena výše; v~ČR naopak zákonné minimum plní roli poslední záchranné sítě pro pracovníky mimo dosah kolektivní smlouvy.

Bodový graf závislosti poměru průměrné ke minimální mzdě (nebo mediánové mzdy k~dolnímu decilu) na \ac{KS} pokrytí ukazuje, jak silné \ac{KV} instituce stlačují vzdálenost mezi minimálním a~mediánovým standardem: v~zemích s~erga omnes je tento poměr nižší --- tedy minimální mzda (nebo kolektivní smluvní tarif) je blíže mediánu.
V~ČR je zákonná minimální mzda relativně vzdálena od mediánu --- přibližně \SI{53}{\percent} mediánu v~roce 2023, zatímco v~DK nebo BE nepotřebují zákonnou minimální mzdu vůbec, protože smluvní minimum je nastaveno kolektivně vysoko.
Tento graf demonstruje, že minimální mzda a~\ac{KS} pokrytí jsou substituty: v~zemích se silným \ac{KV} pokrytím je zákonná minimální mzda méně relevantní (\ac{KS} ji nastavuje výše), zatímco v~zemích s~nízkým pokrytím (ČR) zákonná minimální mzda je posledním záchranným sítem --- ale sítem s~hrubými oky.
Politická implikace: namísto soupeření o~výši zákonné minimální mzdy by měla být prioritou reforma \ac{KV} pokrytí, která by zákonné minimum přirozeně odsunula do role pouhé bezpečnostní sítě pro ty, kteří jsou mimo dosah \ac{KS} --- čímž by se oba nástroje staly komplementárními, nikoliv soupeřícími.

Bodový graf závislosti poměru průměrné ke minimální mzdě (nebo mediánové mzdy k~dolnímu decilu) na \ac{KS} pokrytí ukazuje, jak silné \ac{KV} instituce stlačují vzdálenost mezi minimálním a~mediánovým standardem: v~zemích s~erga omnes je tento poměr nižší --- tedy minimální mzda (nebo kolektivní smluvní tarif) je blíže mediánu.
V~ČR je zákonná minimální mzda relativně vzdálena od mediánu --- přibližně \SI{53}{\percent} mediánu v~roce 2023, zatímco v~DK nebo BE nepotřebují zákonnou minimální mzdu vůbec, protože smluvní minimum je nastaveno kolektivně vysoko.
Tento graf demonstruje, že minimální mzda a~\ac{KS} pokrytí jsou substituty: v~zemích se silným \ac{KV} pokrytím je zákonná minimální mzda méně relevantní (\ac{KS} ji nastavuje výše), zatímco v~zemích s~nízkým pokrytím (ČR) zákonná minimální mzda je posledním záchranným sítem --- ale sítem s~hrubými oky.
Politická implikace: namísto soupeření o~výši zákonné minimální mzdy by měla být prioritou reforma \ac{KV} pokrytí, která by zákonné minimum přirozeně odsunula do role pouhé bezpečnostní sítě pro ty, kteří jsou mimo dosah \ac{KS} --- čímž by se oba nástroje staly komplementárními, nikoliv soupeřícími.

Bodový graf závislosti poměru průměrné ke minimální mzdě (nebo mediánové mzdy k~dolnímu decilu) na \ac{KS} pokrytí ukazuje, jak silné \ac{KV} instituce stlačují vzdálenost mezi minimálním a~mediánovým standardem: v~zemích s~erga omnes je tento poměr nižší --- tedy minimální mzda (nebo kolektivní smluvní tarif) je blíže mediánu.
V~ČR je zákonná minimální mzda relativně vzdálena od mediánu --- přibližně \SI{53}{\percent} mediánu v~roce 2023, zatímco v~DK nebo BE nepotřebují zákonnou minimální mzdu vůbec, protože smluvní minimum je nastaveno kolektivně vysoko.
Tento graf demonstruje, že minimální mzda a~\ac{KS} pokrytí jsou substituty: v~zemích se silným \ac{KV} pokrytím je zákonná minimální mzda méně relevantní (\ac{KS} ji nastavuje výše), zatímco v~zemích s~nízkým pokrytím (ČR) zákonná minimální mzda je posledním záchranným sítem --- ale sítem s~hrubými oky.
Politická implikace: namísto soupeření o~výši zákonné minimální mzdy by měla být prioritou reforma \ac{KV} pokrytí, která by zákonné minimum přirozeně odsunula do role pouhé bezpečnostní sítě pro ty, kteří jsou mimo dosah \ac{KS} --- čímž by se oba nástroje staly komplementárními, nikoliv soupeřícími.

\input{texparts/python/coverage_ratio_scatter}




\paragraph{Predikce vlivu na hospodářství ČR: hypotetické pokrytí~80~\%}\label{par:predikce}

Směrnice Evropského parlamentu a~Rady (EU) 2022/2041 o~přiměřených minimálních mzdách ukládá členským státům s~pokrytím \ac{KS} pod hranicí 80~\% povinnost vypracovat akční plán pro jeho zvýšení.
ČR se svým pokrytím $\approx 31\,\%$ (2023) patří k~zemím s~největším potenciálním nárůstem v~EU --- potřebný nárůst $\Delta x = 80{,}0 - 30{,}7 = +49{,}3$~pp je po Litvě a~Polsku třetím největším v~EU26.

Predikovaná změna výstupních proměnných při tomto nárůstu se dle rov.~\eqref{eq:delta_pred} vypočte jako:
\begin{equation}\label{eq:model_cz}
  \Delta\hat{y}_{\text{CZ}}
    = \hat{\beta}\,(x_1 - x_0)
    = \hat{\beta} \cdot 49{,}3\ \text{pp},
\end{equation}
kde $\hat{\beta}$ jsou panelové odhady dle rov.~\eqref{eq:ols} a~$x_0 = 30{,}7\,\%$, $x_1 = 80{,}0\,\%$.
Výsledky jsou shrnuty v~tab.~\ref{tab:model_cz}.

\begin{table}[htbp]
\centering
\caption{Modelový výpočet dle rov.~\eqref{eq:model_cz}: predikovaná změna ukazatelů ČR
         při nárůstu pokrytí \ac{KS} na~80~\,\% (panelový $\hat{\beta}$, EU26, 2004--2024).}
\label{tab:model_cz}
\begin{tabular}{lrrrr}
\toprule
Ukazatel
  & Aktuální $y_0$
  & $\hat{\beta}$
  & $\Delta\hat{y} = \hat{\beta}\cdot 49{,}3$
  & Predikce $\hat{y}_1$ \\
\midrule
Průměrná prac.~doba (h/týden)    & $39{,}2$ & $-0{,}042$ & $-2{,}1$  & $37{,}1$ \\
HDP/ob.~v~PPS (EU27\,=\,100)     & $92{,}4$ & $+0{,}476$ & $+23{,}5$ & $115{,}9$ \\
Giniho koeficient (0--100)        & $24{,}0$ & $-0{,}037$ & $-1{,}8$  & $22{,}2$ \\
Čistý příjem (EUR~PPS/rok) & \num{19556} & $+123{,}5$ & $+\num{6089}$ & \num{25645} \\
\bottomrule
\end{tabular}
\end{table}

Výsledky je třeba interpretovat jako horní odhad korelovatelné části efektu, nikoliv jako kauzální predikci: $\hat{\beta}$ je odhadnut na průřezu zemí s~různými institucemi, historií a~ekonomickými strukturami a~nezachycuje zpětné vazby (general equilibrium) ani náklady přechodu.
Přesto jsou závěry konzistentní s~tezí práce:
\begin{itemize}
  \item nárůst čistého příjmu o~$+\num{6089}$~EUR~PPS/rok ($+31\,\%$) by ČR přiblížil aktuální úrovni Slovinska nebo Španělska;
  \item zkrácení průměrné pracovní doby o~$2{,}1$~h/týden při vyšším příjmu odpovídá přibližně jednomu dni volna za 14~dní;
  \item pokles Gini o~1,8~bodu odpovídá zhruba pěti letům přirozené konvergence uvnitř EU.
\end{itemize}
Reforma \ac{KV} pokrytí (zavedení mechanismu erga omnes, revize zákona o~kolektivním vyjednávání č.\,2/1991~Sb.~\cite{zakon_kv_1991}) je z~pohledu dat jednou z~nejméně nákladných dostupných politik: nevyžaduje přímé státní výdaje a~generuje distribuční efekty srovnatelné s~rozsáhlými fiskálními programy.
